\documentclass[10pt,twoside]{book}
\usepackage[utf8]{inputenc}
\usepackage[T1]{fontenc}
\usepackage{geometry}
\geometry{margin=1in}
\usepackage{titlesec}
\usepackage{titletoc}
\usepackage{fancyhdr}
\usepackage{xcolor}
\usepackage{tikz}
\usetikzlibrary{shapes.geometric, arrows.meta, positioning}
\usetikzlibrary{shapes,arrows,positioning}
\usepackage{enumitem}
\usepackage{setspace}
\usepackage{sectsty}
\usepackage{needspace}
\usepackage{changepage}
\usepackage{array}
\usepackage{booktabs}
\usepackage[most]{tcolorbox}
\tcbuselibrary{breakable}
\tcbset{breakable}
\usepackage{hyperref}
\hypersetup{colorlinks=true,linkcolor=blue,citecolor=blue,urlcolor=blue}

% Custom boxes
\newtcolorbox{adventurebox}[1][]{colback=gray!10,colframe=black!70,fonttitle=\bfseries,title=#1}
\newtcolorbox{mechanicbox}[1][]{colback=gray!15,colframe=gray!60,fonttitle=\bfseries,title=#1}
\newtcolorbox{npcbox}[1][]{colback=blue!5,colframe=blue!50,fonttitle=\bfseries,title=#1}
\newtcolorbox{timersidebar}[1][]{colback=red!5,colframe=red!70,fonttitle=\bfseries,title=#1}
\newtcolorbox{tipsbox}[1][]{colback=green!5,colframe=green!50,fonttitle=\bfseries,title=#1}
\newtcolorbox{advancementbox}[1][]{colback=purple!5,colframe=purple!50,fonttitle=\bfseries,title=#1}
\newtcolorbox{recapbox}[1][]{colback=yellow!10,colframe=orange!70,fonttitle=\bfseries,title=#1}
\newtcolorbox{cursebox}[1][]{colback=brown!5,colframe=brown!70,fonttitle=\bfseries,title=#1}

% Title formatting
\titleformat{\chapter}{\Huge\bfseries}{\thechapter}{20pt}{}
\titleformat{\section}{\Large\bfseries}{\thesection}{15pt}{}
\titleformat{\subsection}{\large\bfseries}{\thesubsection}{12pt}{}

% Headers
\pagestyle{fancy}
\fancyhf{}
\fancyhead[LE,RO]{\thepage}
\fancyhead[RE]{\nouppercase{\leftmark}}
\fancyhead[LO]{\nouppercase{\rightmark}}

\begin{document}

% Title page
\begin{titlepage}
\centering
\vspace*{2cm}
{\Huge\bfseries Beyond Millhaven\par}
\vspace{1cm}
{\Large The Ford of Three Promises\par}
\vspace{2cm}
{\large A Starter Adventure for Fate's Edge\par}
\vspace{1cm}
{\normalsize \textit{Sequel to ``Blood and Silk''}\par}
\vspace{2cm}
{\normalsize Designed for 4--6 players, Transition from Rookie to Seasoned (30--60 XP)\par}
\vspace{1cm}
{\normalsize Game Master's Guide Included\par}
\vfill
{\normalsize Featuring: Strings • Faction Timers • Mediation • Threshold Negotiation • Advancement\par}
\vspace{0.5cm}
{\normalsize Cover illustration by (placeholder)\par}
\end{titlepage}

\chapter*{How to Run This Adventure}
\addcontentsline{toc}{chapter}{How to Run This Adventure}

\begin{recapbox}[Previously in Blood and Silk]
Your players defeated the Ursillo Gang and saved Millhaven. They earned the trust of the villagers and gained their first Assets (e.g., Village Militia, Elder's Trust). They met \textbf{Tema}, a Vilikari merchant, and may have encountered \textbf{Khatun Sarnai} (Ykrul chieftain) or \textbf{Captain Rusk} (Black Banner mercenary) if you used the optional hooks. Now they carry a \textbf{Letter of Introduction} to a Vilikari lord---their ticket out of exile.

\vspace{0.5em}
\noindent \textit{If your group played the prequel adventure \textbf{The Velvet War of Silkstrand}, the exiles carry more than a letter. They carry the memory of Old Kes, the fence who taught them to keep their word even when it cost everything. That memory may surface here.}
\end{recapbox}

This adventure assumes your players have completed \textit{Blood and Silk} or have reached approximately 30--40 XP. They are no longer desperate exiles---they have proven themselves in Millhaven. Now they face a different kind of challenge: one that cannot be solved with a sword.

\subsection*{Continuity from Millhaven}

The Curse of Acasia does not end at the village boundary. The old Imperial road remembers the broken promise that doomed Millhaven's harvest, and the ford at Valvano was built by the same legionaries who carved the milestones now pointing the wrong way. If your players attracted the Hollow's attention in the previous adventure, that debt is not forgotten---it is merely deferred. You may choose to have a gray figure appear at the ford's edge as the party arrives, or simply let the chill in the air speak for itself.

\section*{The GM's Mindset}

This is a \textbf{negotiation thriller}. Combat is possible, but it is a failure state. The drama comes from competing legitimate claims, not from villainy. Each faction at the ford has a grievance you can sympathize with. Your job is not to decide who is right---it is to make the players choose, and to make that choice matter.

\begin{tipsbox}[GM Mantra]
\textit{Set the DV low. Make the world react honestly. Let the pressure do the work.}
\end{tipsbox}

\section*{Pacing the Three Acts}

\begin{itemize}
\item \textbf{Act I (The Road to the Ford):} 45--60 minutes. Establish the mission, introduce the factions, arrive at the standoff.
\item \textbf{Act II (The Three Promises):} 60--90 minutes. Meet each faction, uncover hidden needs, gather leverage.
\item \textbf{Act III (The Mediation):} 45--60 minutes. The moot, the crossing, the resolution.
\end{itemize}

Total: 2.5--3.5 hours. If players are deeply engaged in negotiation, let it breathe. The timers will keep pressure from stalling.

\section*{Handling the Timers}

The \textbf{Ford Standoff} timer (8 segments) is your primary pacing tool. 

\begin{timersidebar}[Ford Standoff Triggers]
\textbf{Advance +2 segments:}
\begin{itemize}
\item A negotiation fails outright (any faction walks away).
\end{itemize}
\textbf{Advance +1 segment:}
\begin{itemize}
\item A partial success creates new tension (a faction learns of the salt cache).
\item The PCs take too long to act (two scenes without progress).
\end{itemize}
\textbf{Retreat -1 segment:}
\begin{itemize}
\item The PCs broker a partial agreement (two factions agree on one point).
\item The PCs uncover a hidden need and address it.
\item The PCs propose a creative solution that genuinely surprises you.
\end{itemize}
\end{timersidebar}

If the timer fills before a resolution, the standoff collapses. Roll once on the table below to determine who strikes first---then run a desperate escape, not a full battle.

\begin{center}
\begin{tabular}{@{}l l@{}}
\toprule
\textbf{d6} & \textbf{First Striker} \\
\midrule
1--2 & Ykrul riders charge the bridge (honor insult). \\
3--4 & Aeler engineers seal the bridge permanently (collapse risk). \\
5--6 & Black Banner mercenaries loot both sides (Rusk's payday). \\
\bottomrule
\end{tabular}
\end{center}

\section*{Advanced Tips for Tier II Play}

Your players have survived Millhaven. They understand Boons, Timers, and the basic dice loop. Now you can introduce:

\begin{itemize}
\item \textbf{Strings as currency.} Remind players they can offer their existing Strings as leverage. A Ykrul safe-conduct from a previous adventure might sway Sarnai.
\item \textbf{Bonds in negotiation.} When a PC invokes a Bond with another PC during a social roll, grant +1 die. The memory of shared struggle is leverage.
\item \textbf{Debt as a resource.} Rusk's price is a \textbf{Debt timer} (4 segments). The PCs can pay it now (sacrifice a String or asset) or accept the debt and have it come due later.
\end{itemize}

\begin{mechanicbox}[Tier II Transition]
This adventure is designed to bridge Rookie (30--40 XP) and Seasoned (41--90 XP). By the end, each PC should have:
\begin{itemize}
\item Earned 8--12 XP (see Chapter: Advancement)
\item Gained at least one new String
\item Made a meaningful choice about their character's relationship to power
\end{itemize}
\end{mechanicbox}

\section*{If the Players Get Stuck}

The most common stall is analysis paralysis---too many options, no clear path. If the negotiation stalls for more than ten minutes:

\begin{enumerate}
\item \textbf{Advance the timer.} \textit{``The Ykrul grow restless. A rider spurs her horse in a tight circle.''} (Advance Ford Standoff +1)
\item \textbf{Have an NPC make an offer.} \textit{``Tema pulls you aside. 'I can pay Rusk with the salt cache---but I'll need your word you'll protect the seeds. Old Kes spoke well of your people. He said the Silk Coin kept their word. I knew him before the Court took him. He was a good man — broken, but good.' ''}
\item \textbf{Introduce Morvain early.} The Ecktorian prefect is a catalyst; his arrival forces factions to choose sides.
\item \textbf{Ask a leading question.} \textit{``Which faction seems most likely to compromise? What are you willing to offer them?''}
\end{enumerate}

\section*{When to Roll (and When Not To)}

Not every conversation needs a roll. If a PC asks a straightforward question and the NPC has no reason to lie, just answer. Reserve rolls for:

\begin{itemize}
\item Persuading a faction to change a stated position
\item Detecting hidden needs (Insight)
\item Hiding the salt cache (Deception)
\item Interpreting ancient charters (Lore)
\end{itemize}

When a player rolls, state the Position and DV before the dice hit the table. \textit{``This is Controlled, DV 3. Sarnai respects directness.''}

\section*{Safety and Consent}

This adventure deals with honor, shame, and the weight of grudges. Some players may have strong reactions to themes of blood-debt, public apology, or coercive negotiation. Use Lines and Veils as your table requires. The X-card is always available.

\chapter*{Introduction}
\addcontentsline{toc}{chapter}{Introduction}

The road from Millhaven leads west, toward the passes that separate Acasia from the Ecktorian frontier. Beyond the garrison town of Ironbound, the mountains give way to the Violet Steppe. And at the edge of the steppe, where the last stone bridge crosses the River Adurn, stands the Valvano Ford.

Once the seat of a fallen ducal house, now a contested crossing where Ykrul herders, Aeler toll-keepers, and Vilikari merchants stare at one another across a span of gray stone.

You are no longer the desperate exiles who stumbled into Millhaven. You have proven yourselves. But the road does not care about your reputation. The road only cares about the price of crossing.

The caravan master Sorian Tabb offers you honest work: escort a cargo of grain seed to the monastery of St. Arilan in the upper Karasus basin. The pay is modest. The real prize is a Letter of Introduction---a chance to leave Acasia behind and start anew among the Vilikari river lords.

But the ford is closed. A silence furlong seals the east bank. A keystone clause locks the bridge. And a Black Banner mercenary counts coin while three factions refuse to move.

No one can break the deadlock. No one, perhaps, except those who have already learned that every choice has a price---and that some debts are worth incurring.

\subsection*{The Weight of Acasia}

The Curse follows the old roads. The milestones west of Millhaven are chiseled in two languages: Utaran, worn smooth by a thousand rains, and a newer script that no one can read. The locals say the stones forget your name if you walk past them at dusk. They leave offerings of salt and silence. The party may remember the Hollow's attention from the previous adventure; if the timer was not cleared, a gray figure now walks a mile behind them, never closer, never farther. The ford is a threshold, and thresholds attract attention.

\section*{Design Philosophy}

This adventure is a negotiation thriller. Combat is possible but treated as a failure of diplomacy. The themes---legitimacy, debt, silence, and the weight of oaths---emerge through play, not exposition. Each faction has a legitimate grievance. No one is evil. The players' choices create the resolution.

\section*{What You Need}

\begin{itemize}
\item The Fate's Edge Core Rules
\item The travel deck (or substitute card draws)
\item A handful of tokens for tracking Strings and Timers
\item (Optional) The pre-generated characters from \textit{Blood and Silk}
\end{itemize}

\chapter{The Valvano Ford}

\section{Geography}

The Valvano Ford lies at the western edge of the Acasian marches, where the last Aelerian foothills give way to the Ecktorian frontier. A single stone bridge---Aeler-built, Vilikari-chartered---spans the River Adurn. Three routes converge here:

\begin{itemize}
\item \textbf{The Western Road (from Acasia):} Through the Aelerian foothills, past the garrison town of Ironbound.
\item \textbf{The Steppe Track (north-east):} Ykrul seasonal pastures along the Violet Steppe's edge.
\item \textbf{The River Path (south-east):} Along the Karasus toward Vilikari riverlands.
\end{itemize}

Control of the ford is control of access. The Aeler built the bridge and hold the maintenance charter. The Ykrul claim the right of undisturbed passage for their herds---a right older than the charter. The Vilikari need the route open for trade with Acasia and the west. And the Ecktorian prefect at Ironbound watches from a distance, ready to declare the whole crossing an imperial road.

\begin{mechanicbox}[Geography as Leverage]
The ford's position between Acasia, Ecktoria, and the Ykrul steppe means every faction has a different legal claim. Players who pay attention to geography can argue: ``The Ecktorians have no authority here---this is Acasian soil, and Acasia recognizes the Vilikari charter.'' (Presence + Lore, DV 4, grants +1 die to negotiations with Voss or Tema.)
\end{mechanicbox}

\section{The Standoff}

When the party arrives, they find a tense silence. On the east bank, Ykrul riders have raised a \textit{silence furlong}---a ritual border of rope and bone. No speech, no crossing. They demand the Aeler remove iron spikes from the bridge that lame horses, and the Vilikari pay a blood-price for a herder killed last season.

On the west bank, Aeler engineers have locked the bridge's winch with a rune-sealed padlock. Factor-Emissary Voss cites a charter granting his guild the right to collect tolls. He will open the bridge only if the Ykrul recognize the charter and the Vilikari stop smuggling salt by the lower ford.

Between them, the Vilikari factor Tema and a Black Banner captain named Rusk argue over a ledger. Tema's caravan carries grain seed for a famine-relief monastery---and a hidden cache of untaxed salt. Rusk cares only about payment.

\subsection*{Faction Cheat Sheet}

\begin{center}
\begin{tabular}{@{}l l l p{3cm}@{}}
\toprule
\textbf{Faction} & \textbf{Leader} & \textbf{Public Demand} & \textbf{Hidden Need} \\
\midrule
Ykrul Gray Ash & Khatun Sarnai & Remove iron spikes; blood-price & Restore clan honor (public apology) \\
Aeler Guild & Factor Voss & Recognize charter; stop salt smuggling & Show revenue (face-saving audit) \\
Vilikari & Tema (merchant) & Open ford; protect cargo & Hide untaxed salt cache \\
Black Banner & Captain Rusk & Get paid & Any payment (no loyalty) \\
Ecktorian (wild card) & Prefect Morvain & Impose imperial tariff & Legitimacy (bluffing) \\
\bottomrule
\end{tabular}
\end{center}

\chapter{Act I: The Road to the Ford}

\section{Opening Scene: Sorian's Offer}

The party meets Sorian Tabb in a waystation called the Muddy Stilt, two days west of Millhaven. He is a Vilikari caravan master with a broken nose and honest eyes. The inn's hearth is low, and the rain drums on the roof. A single candle gutter on the table between them.

\begin{adventurebox}[Sorian Tabb]
\textit{``I've run this route a dozen times. Never seen the ford locked like this. The Ykrul say they'll wait till the snows; the Aeler say they'll wait till the charter expires; the Vilikari say they'll hire more swords. Me? I just need these seeds to the monastery before the first frost.''}

He offers each PC a letter of introduction to a Vilikari lord---a chance to start over. The pay is three silver marks each, plus a share of any trade negotiated.
\end{adventurebox}

The players may ask questions. Sorian knows:

\begin{itemize}
\item The Ykrul clan is the Gray Ash, led by Khatun Sarnai. Her nephew was killed by a Vilikari caravan that fled without paying wergild.
\item The Aeler factor is Voss, a mid-level bureaucrat under audit by the high council. He cannot afford to lose face.
\item The Black Banner captain is Rusk, a pragmatist. He will side with whoever pays his company's arrears.
\item There is a fourth interested party: Prefect Morvain from Fort Ironbound. He has no authority at the ford, but he carries papers that could create one.
\end{itemize}

If a PC asks Sorian about the Curse, he shrugs. \textit{``The old road is hungry. Leave an offering at the first milestone. A coin, a hair, a breath. The stones remember who was polite.''}

\begin{tipsbox}[Establishing Stakes]
Before the first scene, ask each player: \textit{``What does your character want more than coin?''} Write the answers down. You will use them to tempt hard choices in Act III.
\end{tipsbox}

\section{Travel Leg: The Western Road}

If the GM wishes to include a short journey, draw one card from the Acasia or Ecktorian travel deck. Suggested hazards:

\begin{itemize}
\item \textit{Ecktorian Inspection} (Club): A patrol from Ironbound demands papers. The party must bluff, bribe, or fight. The patrol leader has a glass eye that weeps rust; locals say he sold his real eye to a debt-collector.
\item \textit{Ykrul Scouts} (Club): Riders circle the caravan at dusk, testing their nerve. A successful Presence + Command (DV 3) sends them away respectful. One rider wears a braid of white hair---a token of mourning. Do not ask about it.
\item \textit{The Hollow's Breath} (Club, if Hollow Attention timer advanced in previous adventure): A gray fog rolls in. The milestones seem to shift. A character who rolls a 1 on any navigation test hears their own name whispered from the wrong direction. The GM gains 2 SB.
\end{itemize}

Success on the travel leg grants +1 die to future negotiations with the faction the PCs impressed (e.g., bribing the Ecktorian patrol earns a +1 die with Morvain).

\section{Arrival at the Ford}

Describe the scene:

\begin{quote}
\textit{The river is slate-gray under a low sky. On the far bank, a rope strung between two cairns bears bleached horse-skulls---the silence furlong. Behind it, Ykrul riders sit motionless on shaggy ponies. On this side, the stone bridge rises, its winch sealed with a brass lock etched in Aeler runes. A Vilikari wagon train huddles under canvas, and a scarred man in a black surcoat leans on a two-handed sword.}
\end{quote}

Sorian introduces the party as neutral mediators. \textit{``You have no flag,'' he says, ``so you can talk to all of them without making an enemy.''}

\begin{mechanicbox}[Reading the Scene]
A successful Wits + Notice (DV 3) reveals:
\begin{itemize}
\item The Ykrul riders are restless---their horses keep shifting weight. \textit{(Sarnai is under pressure from her clan.)}
\item The Aeler lock has fresh tool marks. \textit{(Voss has been tampering with his own seal, testing it.)}
\item Rusk's sword is notched. \textit{(He has seen recent combat---he is not just posturing.)}
\item A ninth skull on the silence furlong is smaller than the others. It is a child's. The Ykrul do not speak of it.
\end{itemize}
Each observation grants +1 die to the first social roll against that faction.
\end{mechanicbox}

\subsection*{The Curse at the Ford}

The bridge has nine arches. The locals count only eight. The ninth arch is there, but they will not look at it. Elder Thorne once told a story: a legionary swore on the ninth arch that he would return to pay his toll. He never did. The stone remembers. If a character deliberately crosses the ninth arch, the GM gains 3 SB and the Hollow's attention timer advances by 2. The water beneath that arch is always still, even when the river runs high.

\newpage
\chapter*{Appendix: The Road Remembers}
\addcontentsline{toc}{chapter}{Appendix: The Road Remembers}
\markboth{APPENDIX: THE ROAD REMEMBERS}{APPENDIX: THE ROAD REMEMBERS}

\section*{A Letter from Mareth}

If the players met the exiled scholar in \textit{Blood and Silk}, they may find a folded note tucked into Sorian's satchel when they reach the Muddy Stilt. The paper is coarse, the ink smudged, but the handwriting is unmistakable.

\begin{quote}\itshape
To those who walked the broken road with me:

The ford you seek is a threshold. Thresholds listen. I have learned that the Traveler is not the only patron who watches the crossing places. Another presence lingers there---older, colder, patient. It does not speak. It records. Be careful what you promise, and more careful what you hear.

I have enclosed a few new verses. The road has taught me that silence is also a kind of speech.

\textit{Walk true, but not alone.}

Mareth
\end{quote}

If the players never met Mareth, the letter arrives anyway, addressed to ``the exiles of Millhaven.'' A stranger saw them on the road and thought they might need guidance. The GM may treat this as a minor String: \textit{An Unknown Friend}.

\section*{Deepening the Bond with the Traveler}

Characters who learned low rites from Mareth in the first adventure may now deepen their connection to the Traveler. The GM may offer the following options during downtime or as rewards for roleplaying devotion to the road.

\subsection*{New Low Rites of the Traveler}

These rites can be learned from the letter (no XP cost, but require a downtime scene to practice).

\begin{adventurebox}{Rite of the Unspoken Promise (Low Rite)}
\textbf{Cost:} 1 action, mark +1 Obligation. \\
\textbf{Effect:} For the rest of the scene, you may make a promise without speaking it aloud. The Traveler witnesses the intent, not the words. Anyone who tries to magically compel you to reveal the promise suffers –1 die on the roll. \\
\textbf{Push:} Mark 1 Fatigue to extend the effect to one ally. \\
\textit{Mareth's note: ``Silence is a shield. The ford respects those who do not waste words.''}
\end{adventurebox}

\begin{adventurebox}{Rite of the Forked Sign (Low Rite)}
\textbf{Cost:} 1 action, mark +1 Obligation. \\
\textbf{Effect:} You lay a small stone, twig, or coin at a crossroads or threshold. For the next hour, any creature that crosses that threshold and intends to deceive must roll Resolve (DV 3) or be compelled to speak one true word aloud (GM chooses the word based on context). The sign is obvious to anyone looking for it; stealthy crossing may bypass the effect. \\
\textbf{Push:} Mark +1 Obligation to make the sign last until dawn. \\
\textit{Mareth's note: ``The Aeler use this to catch toll-cheats. The Ykrul use it to test oaths. Use it wisely.''}
\end{adventurebox}

\begin{adventurebox}{Rite of the Weary Step (Low Rite)}
\textbf{Cost:} 1 action, mark +1 Obligation. \\
\textbf{Effect:} You touch a willing creature (or yourself) and murmur a word of rest. The creature ignores the first point of Fatigue from travel or exhaustion during the next leg of a journey. This does not stack with similar effects. \\
\textbf{Push:} Mark 1 Fatigue to extend the effect to two additional creatures. \\
\textit{Mareth's note: ``The road takes more than it gives. Sometimes you must ask it to wait.''}
\end{adventurebox}

\subsection*{The Traveler's Second Gift (Imbuement Upgrade)}

A Runekeeper who has completed at least one adventure (such as \textit{Blood and Silk}) may unlock an enhanced version of the Patron's Gift. This requires a downtime scene of meditation at a crossroads and costs 2 XP (or may be granted as a story reward).

\begin{adventurebox}{Gift of the Unmarked Way (Imbuement, Upgraded)}
\textbf{Cost:} 1 action, mark +1 Obligation. \\
\textbf{Effect:} Touch a held item (weapon, tool, or cloak). For the rest of the scene, the item grants:
\begin{itemize}
\item +2 dice to Navigation, Survival, or Escape rolls (instead of +1).
\item The bearer may ignore one environmental penalty (difficult terrain, fog, darkness) entirely, once per scene.
\end{itemize}
\textit{``The road knows your name now. It will not hide from you.''}
\end{adventurebox}

\section*{A Second Patron: The Witness at the Ford}

Mareth's letter hints at another presence: \textbf{The Witness}, the patron of truth, memory, and inconvenient revelation. At the Valvano Ford, the Witness is strongest near the bridge's keystone, where generations of oaths have been sworn and broken. A character who wishes to learn low rites of the Witness may do so by spending a scene in meditation at the ninth arch (or by successfully mediating a dispute that reveals a hidden truth).

\begin{cursebox}{The Witness's Attention}
Unlike the Traveler, the Witness does not grant Obligation—it grants \textbf{Recognition}. When you invoke a Rite of the Witness, mark +1 Recognition instead of Obligation. Your Recognition Capacity equals your Spirit + Presence. Exceeding capacity does not cause Fatigue; instead, the Witness begins to reveal truths you would rather not know (GM gains 1 SB per excess segment, usable for complications that expose secrets or force inconvenient confessions).
\end{cursebox}

\subsection*{Low Rites of the Witness}

\begin{adventurebox}{Rite of the Lingering Trace (Low Rite)}
\textbf{Cost:} 1 action, mark +1 Recognition. \\
\textbf{Effect:} Choose a trace (a footprint, a broken lock, a discarded object). You may ask the GM one yes/no question about the last creature to touch that trace. The GM must answer truthfully. \\
\textbf{Push:} Mark 1 Fatigue to ask a second question. \\
\textit{``The ford keeps a ledger of every foot that crossed it. Read it carefully.''}
\end{adventurebox}

\begin{adventurebox}{Rite of the Uncomfortable Question (Low Rite)}
\textbf{Cost:} 1 action, mark +1 Recognition. \\
\textbf{Effect:} Ask a single NPC one direct question about a promise they have made or broken. They must either answer truthfully (the GM gives a true answer, though it may be cryptic) or mark 1 Fatigue and gain a \textit{Guilty Tell} ( –1 die to all social rolls until they confess). \\
\textbf{Push:} Mark 1 Story Beat (Hearts) to force the question to be answered in front of witnesses. \\
\textit{``The Ykrul respect those who speak plainly. The Aeler fear those who listen too well.''}
\end{adventurebox}

\begin{adventurebox}{Rite of the Shared Burden (Low Rite)}
\textbf{Cost:} 1 action, mark +1 Recognition. \\
\textbf{Effect:} You take a psychological Condition (Guilty, Haunted, Shamed) from a willing ally and suffer it yourself. Your ally clears the Condition. The Condition lasts for the scene; after the scene ends, you may reduce it to Fatigue 1 instead of carrying it forward. \\
\textbf{Push:} Mark 1 Fatigue to take the Condition without rolling. \\
\textit{``The Witness does not forgive debts. It allows them to be transferred.''}
\end{adventurebox}

\section*{Magic in the Mediation}

The ford is a threshold, and thresholds are where magic is most potent. Characters with magical abilities (Runekeepers, Cantors, Invokers, or free casters) may use their arts during the negotiation scenes. Below are specific applications tied to the adventure's factions and challenges.

\subsection*{Using the Traveler's Rites}

\begin{itemize}
\item \textit{Road-Sense} can detect hidden paths around the ford, potentially revealing the smugglers' lower crossing (relevant to Voss's demand about the salt tariff).
\item \textit{Waymark} can be used to signal safe passage to the Ykrul, convincing them that the Aeler will not spike the bridge again (a successful roll grants +1 die to negotiations with Sarnai).
\item \textit{Rite of the Forked Sign} placed at the caravan's camp can expose any agent of Morvain attempting to spy (reveals hidden eavesdroppers).
\end{itemize}

\subsection*{Using the Witness's Rites}

\begin{itemize}
\item \textit{Rite of the Lingering Trace} on the Aeler lock reveals who last tampered with it (Voss himself, testing his own seal). This knowledge gives the PCs leverage: ``You don't want your auditors to know you've been testing the lock, do you?''
\item \textit{Rite of the Uncomfortable Question} directed at Rusk forces him to admit that his contract is with a middleman, not with the faction he claims to serve. This can reduce his demand (he only needs to cover his expenses, not the full contract).
\item \textit{Rite of the Shared Burden} can be used to take a Ykrul's shame onto oneself, allowing Sarnai to accept a compromise without losing face. The PC carries the shame as a Condition for the rest of the scene, but the Ykrul open the ford.
\end{itemize}

\subsection*{Freeform Magic (Threadweaving) at the Ford}

A character with the Spellcraft talent (or the threadweaving training from Mareth) may attempt improvisational magic. Sample applications:

\begin{itemize}
\item \textit{``I weave a thread of subtle knowing to sense what Voss is hiding.''} (Tags: \texttt{subtle}, \texttt{knowing}) DV 3. Success reveals his audit anxiety; failure alerts his guards.
\item \textit{``I tie a knot in the thread of the Ykrul's anger---just for a moment, just enough to speak.''} (Tags: \texttt{entangling}, \texttt{subtle}) DV 3. On success, Sarnai hesitates for one exchange, giving the PCs a chance to speak without interruption. On a miss, she interprets the magic as an insult (advance Ford Standoff +1).
\item \textit{``I pull the thread of the bridge's memory. Show me who built it.''} (Tags: \texttt{knowing}) DV 2. Success reveals that the Aeler engineer who carved the charter's loophole was a woman named Valeriana. A PC with Lore may recall her name from an old text, granting +1 die to negotiations with Voss (he does not want his guild's dirty laundry aired).
\end{itemize}

\begin{tipsbox}[Magic and the Hollow]
Any magical roll that results in a 1 (Story Beat) may attract the Hollow's attention at the ford. The GM may spend the SB to cause a minor manifestation: a shadow that does not belong to any PC, a bell that rings from the riverbed, or a gray figure glimpsed on the far bank. These should be atmospheric, not punishing. The ford remembers all debts, magical and mundane.
\end{tipsbox}

\section*{Cantors and Songs of the Ford}

A Cantor who learned songs from Mareth may adapt them to the ford's acoustics. The stone bridge carries sound strangely; a whispered word can be heard on the far bank, but a shouted oath vanishes into the river.

\subsection*{New Song: The Bridge's Witness}

\textbf{Requirements:} Performance 2+, Lore 2+, must have learned \textit{Rite of the Uncomfortable Question} (or have the Witness's attention). \\
\textbf{Cost:} 1 action, Performance + Lore roll (DV 4). \\
\textbf{Effect:} For the rest of the scene, any lie spoken within earshot of the singer causes the singer's voice to crack (they may mark 1 Fatigue to keep singing, but the lie is still audible to all present). Additionally, the singer may, once per scene, ask one NPC a yes/no question; the NPC must answer truthfully or the song will repeat the lie back to them in the singer's voice. \\
\textbf{Push:} Mark +1 Corruption (Cantor) or +1 Recognition (if using Witness's rules) to extend the effect to all lies spoken within the ford for the next hour.

\begin{adventurebox}{Cantor's Burden at the Ford}
The bridge's ninth arch is a resonance point. A Cantor who sings from the ninth arch may treat any Performance + Lore roll as having +2 dice, but each 1 rolled generates 2 SB instead of 1. The Hollow is listening. The GM may spend these SB to have the song attract the attention of a gray figure from the river.
\end{adventurebox}

\section*{Invokers and the Ford's Keystone}

An Invoker who carries the Traveler's compass (or a symbol of the Witness) may attempt to \textit{Crack the Seal} on the Aeler lock without negotiating. This is risky but dramatic.

\begin{adventurebox}{Crack the Seal (Invoker)}
\textbf{Cost:} Spend 1 action, set the Symbol to \textit{Compromised}, mark +2 Obligation (or +2 Recognition). \\
\textbf{Effect:} Resolve the Aeler lock's opening as if you had the key. The bridge winch turns, the chains groan, and the ford opens. However, the GM gains 2 SB immediately and may spend them at any point during the crossing to introduce a complication (e.g., a plank breaks, a Ykrul rider takes offense, Rusk demands payment \textit{now}). \\
\textit{``The keystone remembers every hand that forced it. Yours will be added to the list.''}
\end{adventurebox}

\section*{Advancing a Magical Character}

Characters who used magic in \textit{Blood and Silk} may now have enough experience to deepen their path. The GM may offer the following advancements as rewards for roleplaying or XP purchases.

\subsection*{Runekeeper (Traveler) -- Tier II Advancement}

\begin{itemize}
\item \textbf{New Rite:} \textit{Rite of the Unspoken Promise} (see above).
\item \textbf{Imbuement Upgrade:} \textit{Gift of the Unmarked Way} (2 XP).
\item \textbf{Talent:} \textit{Witness to the Road} (4 XP) -- Once per session, you may declare that you have passed this way before. Gain +1 die to any Navigation or Survival roll in a location you have visited at any point in the campaign (the road remembers you).
\end{itemize}

\subsection*{Cantor (Traveler) -- Tier II Advancement}

\begin{itemize}
\item \textbf{New Song:} \textit{The Bridge's Witness} (requires downtime study, 4 XP).
\item \textbf{Talent:} \textit{Echo of the Ford} (6 XP) -- When you sing a Traveler song at a threshold (bridge, gate, crossroads), you may extend its duration to the entire scene without additional cost. The GM gains 1 SB (the threshold remembers your voice).
\end{itemize}

\subsection*{Invoker (Traveler or Witness) -- Tier II Advancement}

\begin{itemize}
\item \textbf{New Symbol:} \textit{Keystone Shard} (4 XP) -- A chip of stone from the Valvano Ford's ninth arch. When you display this symbol openly, reduce Obligation cost for Traveler rites by 2 (minimum 0) and reduce Recognition cost for Witness rites by 1. The shard is \textit{Compromised} after three uses; it can be restored by a downtime ritual at the ford (DV 4 Lore or Craft).
\item \textbf{Talent:} \textit{Threshold Memory} (8 XP) -- Once per session, you may touch a threshold (door, bridge, gate) and ask the GM one question about an oath or promise made there in the past. The GM must answer truthfully, but the answer may be a riddle.
\end{itemize}

\subsection*{Free Caster (Threadweaver) -- Tier II Advancement}

\begin{itemize}
\item \textbf{New Tags:} \textit{binding} (restrains a promise, not a person---can prevent an NPC from breaking a specific oath for one scene; +2 DV), \textit{recording} (captures a spoken conversation onto a scrap of parchment; the parchment crumbles after being read aloud once; +1 DV).
\item \textbf{Talent:} \textit{Resonant Thread} (6 XP) -- When you successfully cast a threadweaving spell at a threshold (bridge, gate, ford), you gain 1 Boon and the GM gains 1 SB (the threshold is a witness).
\end{itemize}

\section*{A Final Word from Mareth (GM Eyes Only)}

The letter is real, but Mareth is not at the ford. They never intended to be. They are following their own road, somewhere in the mountains, seeking a way to unbind a promise they made years ago. If the players wish to seek them out, that is a seed for a future adventure. For now, the letter serves as a reminder: magic is a tool, not a solution. The ford will open when the promises are balanced, not when a spell is cast.

The Traveler watches. The Witness records. The Hollow counts.

\vspace{1cm}
\begin{center}
\emph{``The road does not end. It only turns.''}
\end{center}

\end{document}